\section{Generators and Relations for an Elementary System}

\paragraph{\textbf{6.1.}}
\SGARefTarget{sga3:xx:6.1}
Let \(S\) be a scheme and let \((G,T,\alpha)\) be an elementary
\(S\)-system.  Let
\[
 X\in W(\mathfrak g^\alpha)^\times(S),
 \qquad u=\exp(X).
\]
As observed in \SGARefLink{sga3:xx:3.8}{3.8}, the element \(w=w_\alpha(X)\) satisfies, in
particular,
\[
 (wu)^3=e.
\]
\footnote{The editors added the sentence that follows.}
Let \(s_\alpha\) denote the automorphism of \(T\) induced by
\(\intaut(w)\).  By \SGARefLink{sga3:xx:3.1}{Theorem 3.1}\textup{\SGARefLink{sga3:xx:3.1:item:iii}{(iii)}}, for every \(S'\to S\)
and every \(t\in T(S')\),
\[
 s_\alpha(t)=\intaut(w)(t)
 =t\,\alpha^\ast\bigl(\alpha(t)^{-1}\bigr).
\]

\paragraph{\textbf{Theorem 6.2.}}
\SGARefTarget{sga3:xx:6.2}
\emph{Let \(H\) be a sheaf of groups over \(S\) for the fppf
topology.  Let}
\[
 f_T\colon T\longrightarrow H,
 \qquad
 f_\alpha\colon U_\alpha\longrightarrow H
\]
\emph{be group morphisms, and let \(h\in H(S)\).  There exists a
\textup{(necessarily unique)} group morphism}
\[
 f\colon G\longrightarrow H
\]
\emph{extending \(f_T\) and \(f_\alpha\) and satisfying \(f(w)=h\)
if and only if the following conditions hold:}

\emph{\textup{(i)} For every \(S'\to S\),
\(t\in T(S')\), and \(x\in U_\alpha(S')\),}
\begin{equation}
\tag{1}\label{sga3:xx:theorem:6-2:equation:1}
 f_T(t)f_\alpha(x)f_T(t)^{-1}
 =f_\alpha(txt^{-1})
 =f_\alpha\bigl(x^{\alpha(t)}\bigr).
\end{equation}
\emph{Equivalently, \(f_T\) and \(f_\alpha\) extend to a group
morphism \(T\cdot U_\alpha\to H\).}

\emph{\textup{(ii)} For every \(S'\to S\) and \(t\in T(S')\),}
\begin{equation}
\tag{2}\label{sga3:xx:theorem:6-2:equation:2}
 h f_T(t)h^{-1}
 =f_T(s_\alpha(t))
 =f_T\!\left(t\,\alpha^\ast(\alpha(t)^{-1})\right).
\end{equation}

\emph{\textup{(iii)} The following two relations hold in \(H(S)\):}
\begin{equation}
\tag{3}\label{sga3:xx:theorem:6-2:equation:3}
 h^2=f_T\bigl(\alpha^\ast(-1)\bigr),
\end{equation}
\begin{equation}
\tag{4}\label{sga3:xx:theorem:6-2:equation:4}
 \bigl(hf_\alpha(u)\bigr)^3=e.
\end{equation}

\smallskip
\noindent\emph{Proof.}
Write the laws on \(U_\alpha\) and \(U_{-\alpha}\) additively, while
writing their vector-space structures multiplicatively.  If \(f\)
has the asserted properties, then for
\(y\in U_{-\alpha}(S')\) one necessarily has
\[
 f(y)
 =f(w^{-1}wyw^{-1}w)
 =h f_\alpha(w^{-1}yw)h^{-1}.
\]
Define a morphism \(f_{-\alpha}\colon U_{-\alpha}\to H\) by
\begin{equation}
\tag{\(\ast_1\)}\label{sga3:xx:theorem:6-2:equation:ast-1}
 f_{-\alpha}(y)=h f_\alpha(w^{-1}yw)h^{-1}.
\end{equation}
It is a group morphism.  On the big cell \(\Omega\), the morphism
\(f\) is forced to satisfy
\[
 f(ytx)=f_{-\alpha}(y)f_T(t)f_\alpha(x).
\]
This proves uniqueness.  The conditions in the statement are
manifestly necessary; let us prove that they are sufficient.

Relation \textup{\SGARefLink{sga3:xx:theorem:6-2:equation:4}{(4)}} gives
\[
 h f_\alpha(u)h^{-1}h^2
 =
 f_\alpha(-u)h^{-1}f_\alpha(-u).
\]
By \textup{\SGARefLink{sga3:xx:theorem:6-2:equation:3}{(3)}} and \textup{\SGARefLink{sga3:xx:theorem:6-2:equation:1}{(1)}},
\(h^2=f_T(\alpha^\ast(-1))\) commutes with \(f_\alpha(-u)\), and
hence
\[
 h f_\alpha(u)h^{-1}
 =
 f_\alpha(-u)h f_\alpha(-u).
\]
By definition,
\[
 h f_\alpha(u)h^{-1}=f_{-\alpha}(wuw^{-1}),
\]
whereas \SGARefLink{sga3:xx:3.7}{3.7} gives
\begin{equation}
\tag{\(\ast_2\)}\label{sga3:xx:theorem:6-2:equation:ast-2}
 wuw^{-1}=-\widetilde u,
\end{equation}
with \(\widetilde u\) the element paired with \(u\).  Therefore
\begin{equation}
\tag{\(\ast_3\)}\label{sga3:xx:theorem:6-2:equation:ast-3}
 f_{-\alpha}(-\widetilde u)
 =
 f_\alpha(-u)h f_\alpha(-u).
\end{equation}

Let \(t\) be a section of \(T\) over a variable \(S'\to S\).
Apply \(\intaut(f_T(t))\) to this identity.  On the left one
obtains\footnote{The editors corrected the calculation that follows.}
\[
\begin{aligned}
 f_T(t)f_{-\alpha}(-\widetilde u)f_T(t)^{-1}
 &=
 f_T(t)h f_\alpha(u)h^{-1}f_T(t)^{-1}\\
 &=
 h\bigl(h^{-1}f_T(t)h\bigr)f_\alpha(u)
   \bigl(h^{-1}f_T(t)^{-1}h\bigr)h^{-1}\\
 &=
 h f_T(s_\alpha(t))f_\alpha(u)
   f_T(s_\alpha(t))^{-1}h^{-1}\\
 &=
 h f_\alpha\bigl(\alpha(s_\alpha(t))u\bigr)h^{-1}\\
 &=
 h f_\alpha\bigl(\alpha(t)^{-1}u\bigr)h^{-1}\\
 &=
 f_{-\alpha}\bigl(-\alpha(t)^{-1}\widetilde u\bigr).
\end{aligned}
\]
Here we used \textup{\SGARefLink{sga3:xx:theorem:6-2:equation:2}{(2)}}, \textup{\SGARefLink{sga3:xx:theorem:6-2:equation:1}{(1)}},
\(s_\alpha(t)=t\alpha^\ast(\alpha(t)^{-1})\), and
\(\alpha\circ\alpha^\ast=2\).

The right-hand side of \((\ast_3)\), after conjugation, becomes
\[
 f_\alpha(-\alpha(t)u)\,
 f_T(t)h f_T(t)^{-1}h^{-1}\,
 h f_\alpha(-\alpha(t)u).
\]
Since
\[
 h f_T(t)^{-1}h^{-1}
 =f_T(s_\alpha(t^{-1}))
  =f_T\bigl(t^{-1}\,\alpha^\ast(\alpha(t))\bigr),
\]
this equals
\[
 f_\alpha(-\alpha(t)u)\,
 f_T\bigl(\alpha^\ast(\alpha(t))\bigr)\,
 h f_\alpha(-\alpha(t)u).
\]
Comparison gives
\[
 f_{-\alpha}\bigl(-\alpha(t)^{-1}\widetilde u\bigr)
 =
 f_\alpha(-\alpha(t)u)\,
 f_T\bigl(\alpha^\ast(\alpha(t))\bigr)\,
 h f_\alpha(-\alpha(t)u).
\]
Since \(\alpha\colon T\to\mathbb G_{m,S}\) is faithfully flat and
\(H\) is a separated presheaf, \(-\alpha(t)^{-1}\) may be replaced
by an arbitrary section of \(\mathbb G_{m,S}\).  Thus:

\paragraph{\textbf{Lemma 6.2.1.}}
\SGARefTarget{sga3:xx:6.2.1}
\emph{For every \(S'\to S\) and every
\(z\in\mathbb G_m(S')\),}
\[
 f_{-\alpha}(z\widetilde u)
 =
 f_\alpha(z^{-1}u)\,
 f_T\bigl(\alpha^\ast(-z^{-1})\bigr)\,
 h f_\alpha(z^{-1}u).
\]

